L'Ainhoa està intrigada per saber a quina altura exploten els coets llançats a la revetlla de Sant Joan. Per a poder-ho determinar se situa a una distància de 50 metres del punt on es llancen els coets i enregistra, amb un sonòmetre, un nivell d'intensitat sonora de 100 decibels en l'explosió d'un coet que no s'ha enlairat.

\begin{apartat}{1,25}
Quina potència sonora emet el coet en el moment de l'explosió? Si l'explosió ha durat 0,03~s, quina energia sonora s'ha alliberat?
\end{apartat}

\begin{apartat}{1,25}
Des de la mateixa distància al punt de llançament, enregistra 90 decibels d'intensitat sonora en el cas d'un coet igual a l'anterior que s'ha enlairat verticalment i ha explotat a certa altura. Calculeu a quina altura ha explotat el coet. Si dos coets idèntics a l'anterior exploten simultàniament a la mateixa altura que abans, quin nivell d'intensitat sonora percebrà l'Ainhoa, si està situada a la mateixa posició d'abans?
\end{apartat}

\nota{Considereu que les ones sonores es propaguen en les tres dimensions de l'espai i la seva energia es distribueix en superfícies esfèriques.}
\dades{%
  $I_0 = 10^{-12}\ \mathrm{W\,m^{-2}}$.\\
  La velocitat del so en l'aire és de 340~m/s.\\
  Superfície esfèrica: $4\pi r^2$.}
